\section{Singular entrances in semialgebraic families}\label{sec:singular-entrances}
The regular-corner calculation of Section~\ref{sec:remote-general} does
not exhaust remote identification. We now describe entrance orders
without a corner parametrization or an injective observation derivative.
The required finiteness is finiteness of the exact parameter fibre,
modulo any already specified finite symmetry, not smoothness of its
representatives.

Let $X\subset\Theta\times K$ be semialgebraic with compact fibres in a
fixed compact set $K$. Let $f_\theta:X_\theta\to\R^M$ and
$A_\theta:X_\theta\to\C[z]$ be continuous semialgebraic maps, the latter
monic of fixed degree. For the probability experiment, $f$ is its
weighted square-root embedding. Let
\[
 p_\theta(\delta)=f_\theta(x_\theta(\delta)),\qquad
 x_\theta(\delta)\longrightarrow x_{0,\theta},\qquad \delta\downarrow0,
\]
be a semialgebraic centre path with $p_\theta$ extending $C^1$ to zero.
All data, including this path, have fixed format. Assume
$f_\theta^{-1}(p_\theta(0))$ is finite. After a finite semialgebraic
partition its elements can be listed by continuous semialgebraic
sections $x_{0,\theta},\ldots,x_{N,\theta}$; subsequent refinements
will be understood. Choose pairwise disjoint small closed relative
neighbourhoods $N_{j,\theta}$ of these points, and put
\begin{equation}\label{eq:singular-distance}
 d_j(\theta,\delta)=\min_{x\in N_{j,\theta}}
              \|f_\theta(x)-p_\theta(\delta)\|.
\end{equation}
The neighbourhoods can be chosen semialgebraically, for example using
closed balls of radius smaller than a quarter of the least separation of
the listed fibre points.

\begin{theorem}[Finite entrance orders and exact membership]\label{thm:singular-entrances}
There is a finite Nash stratification of the parameter set with the
following properties. For each $j$ either $d_j$ is identically zero for
all sufficiently small positive $\delta$, or there are a rational
$\rho_j\ge1$ and a positive Nash function $a_j$ such that, for every
compact subset $H$ of the stratum,
\begin{equation}\label{eq:singular-power}
 d_j(\theta,\delta)=a_j(\theta)\delta^{\rho_j}
            +O_H(\delta^{\rho_j+\nu_j}),\qquad \nu_j>0.
\end{equation}
The nonzero exponents belong to a finite rational list determined by the
fixed-format presentation. The germ of $d_j$ is independent of the
sufficiently small isolating neighbourhood.

For any specified positive semialgebraic threshold
$t_\theta(\delta)\to0$, a further finite partition decides, for each
$j$, whether $t_\theta-d_j$ is eventually positive, negative, or
identically zero. The remote component is included precisely in the first
and third cases. At every fixed parameter the feasible target root sets
converge in Hausdorff distance to
\begin{equation}\label{eq:entrance-union}
 \bigcup_{j:\ d_j\le t_\theta\ \mathrm{eventually}}
                         \{Z(A_\theta(x_{j,\theta}))\}.
\end{equation}
All the strata, exponents, leading-coefficient graphs and membership
decisions are effectively describable for algebraic input. No
regular-corner assumption is made in these assertions.
\end{theorem}
\begin{proof}
First fix a parameter. The boundary of an isolating neighbourhood
contains no point of the exact fibre. Its observation distance from
$p(0)$ has a positive minimum. The same is true of the complement of
smaller neighbourhoods of all the exact points. Minima over these compact
sets are semialgebraic functions of the parameter. Refinement makes them
continuous and positive, so the resulting separation bounds are uniform
on compact subsets of a stratum. Radii and the finitely many fibre
sections admit the same refinement. Finite semialgebraic choice of the
fibre points follows by cell decomposition of their finite-fibre graph.

The competitor $x_j$ gives $d_j\le\|p(\delta)-p(0)\|=O(\delta)$.
If a sequence of minimizers did not tend to $x_j$, compactness would
give a different point of the exact fibre in $N_j$. This is impossible.
The same argument for the annulus between two smaller isolating
neighbourhoods shows that their minima are equal for all sufficiently
small $\delta$. Their distance germs, not just their leading powers,
are therefore independent of this choice.

The graph of \eqref{eq:singular-distance} is a first-order formula:
there is a point attaining the displayed squared distance, and every
point of $N_j$ has at least that squared distance. Quantifier elimination
makes the graph semialgebraic. An eventually nonzero nonnegative
one-variable semialgebraic function has an expansion
$a\delta^\rho+o(\delta^\rho)$ with $a>0$ and rational $\rho$.
For clarity about finiteness in families, apply
Lemma~\ref{lem:support-order} to a quantifier-free graph formula for
$d_j$. Among the finitely many nonzero specialized polynomial supports,
two lowest-order monomials must balance on a positive graph branch.
Their exponent differences give a finite rational candidate list.
Eventual vanishing is tested separately. The comparability test used in
\eqref{eq:order-decision} selects the actual exponent by a first-order
formula. The inequality $d_j=O(\delta)$ forces $\rho\ge1$.

The limit $a=\lim d_j/\delta^\rho$ also has a first-order graph.
Refine until it is positive and Nash. On such a stratum the error
$|d_j/(a\delta^\rho)-1|$ tends to zero on each fibre.
Proposition~\ref{prop:decay} gives a positive rational power bound and
a positive continuous radius. Taking their bounds and the bounds on $a$
on compact subsets proves \eqref{eq:singular-power}. Thus neither a
metric rate nor uniformity has been deduced from topological triviality.

For the threshold assertion, a one-variable semialgebraic function has
constant sign or is zero on some terminal interval $(0,\eta)$. Each
of these alternatives has a first-order formula with the parameter
free, so a finite partition decides them. The sign may be determined
by a higher Puiseux coefficient when leading terms coincide; no finite
truncation is declared sufficient before the sign test has been made.
If the difference is identically zero, compact attainment in
\eqref{eq:singular-distance} puts a point of $N_j$ in the closed ball.
Thus equality is inclusion, not an assigned side of a critical ray.

The exterior separation already proved forces every small-radius
competitor into one of the $N_j$. Within an included neighbourhood,
all feasible points tend to its unique exact point, while a minimizer
provides at least one feasible point for every sufficiently small
$\delta$. Continuity of the target and of finite-degree roots proves
both Hausdorff inclusions in \eqref{eq:entrance-union}. The centre path
ensures the union is nonempty. Finally every selection, limiting graph
and sign test used above is a finite real-algebraic formula. The
finite-support algorithm and real-root isolation in
Theorem~\ref{thm:effective} prove the stated effectivity.
\end{proof}

This theorem describes singular \emph{entrance orders}, not a
classification of singularities of the source. Its proof uses the
relative semialgebraic machinery for a conclusion not supplied by a
regular tangent-cone projection: entrances tangent to the centre path
may occur at strictly higher rational orders.

\begin{theorem}[Realization of every rational contact order]\label{thm:entrance-realization}
Every rational $\rho\ge1$ occurs as a positive remote entrance exponent
in a compact semialgebraic model with strictly positive probability
observations, a two-point exact fibre and distinct polynomial targets
on its two representatives. For $\rho=n/m>1$ the probability maps can
be chosen polynomial in their respective branch parameters, of degree
at most $n$. Hence there is no universal finite list independent of
format, although Theorem~\ref{thm:singular-entrances} gives a finite
list for every fixed-format family.
\end{theorem}
\begin{proof}
Choose a strictly positive probability vector $P_0$ with at least three
cells and linearly independent zero-sum vectors $A,B$ such that
$\langle A,B\rangle_I=0$, where
$\langle H,H'\rangle_I=\sum_eH_eH'_e/P_{0,e}$.
Let the model be the disjoint union of two sufficiently short closed
intervals, with maps
\[
 P_{\rm base}(s)=P_0+As,\qquad
 P_{\rm rem}(u)=P_0+A u^m+B u^n,\qquad s,u\ge0.
\]
All probabilities are positive after shortening the intervals.
Independence of $A,B$ makes the exact fibre at $P_0$ precisely the two
origins. Give the two intervals different constant monic degree-one
targets. Take the centre path $s=\delta$ and use Hellinger distance.

Put $x=u^m$ and $\rho=n/m>1$. The trial $x=\delta$ gives an entrance
upper bound $O(\delta^\rho)$. Hellinger distance is locally comparable
to Euclidean probability distance, and independence gives
\[
 |x-\delta|+x^\rho\le C\|A(x-\delta)+Bx^\rho\|.
\]
Every minimizer thus has $x=\delta+O(\delta^\rho)$. Write
$x=\delta+v\delta^\rho$; the relevant $v$ are bounded. The square-root
embedding, expanded uniformly at $P_0$, gives
\[
 \delta^{-\rho}h(P_{\rm rem}(u),P_{\rm base}(\delta))
     =\tfrac12\|vA+B\|_I+o(1).
\]
Every bounded $v$ is feasible for small $\delta$ since $\rho>1$.
The minimum of the right side is $\|B\|_I/2>0$ by the chosen
orthogonality. This proves the desired exponent and leading coefficient.
For $\rho=1$, use instead $P_{\rm rem}(u)=P_0+Bu$, with the same
orthogonal $A,B$. Projection onto the nonnegative $B$ ray gives the
positive coefficient $\|A\|_I/2$. These constructions concern the
stated general class of probability models; they do not assert that
every exponent occurs in the cubic binary family.
\end{proof}

\subsection{An algebraic test for the regular regime}
The following criterion derives the corner hypotheses from a local
presentation. It also separates two ways for that regime to fail:
singular feasible geometry and tangency of the centre path to a regular
observation image.

\begin{proposition}[A checkable corner criterion]\label{prop:corner-criterion}
At a remote representative $x_*$ suppose the model is locally
\[
 X=\{x:h(x)=0,\ g_i(x)\ge0\},
\]
with Nash functions. Assume $Dh(x_*)$ has full row rank and the
differentials of the active $g_i$, restricted to $\ker Dh(x_*)$, are
linearly independent. Put
\[
 K_* =\{v:Dh(x_*)v=0,\ Dg_i(x_*)v\ge0\text{ for active }i\}.
\]
Suppose $f$ has a $C^2$ Nash extension near $x_*$, the centre admits
$p(\delta)=p(0)+\delta b+\delta^2c+o(\delta^2)$, and
\begin{equation}\label{eq:algebraic-corner-tests}
 \ker Df(x_*)\cap K_* =\{0\},\qquad
 p'(0)\notin Df(x_*)K_*.
\end{equation}
Then a full Nash corner chart exists, its observation derivative is
cone-coercive, and the entrance has exponent one. The first and second
coefficients are exactly those of
Theorems~\ref{thm:remote-general} and \ref{thm:critical-window}.
For algebraic input all these hypotheses are finite algebraic or
first-order tests.
\end{proposition}
\begin{proof}
The Nash implicit-function theorem first supplies coordinates on $h=0$.
The independent active functions can then be used as the first
coordinates on this manifold, completed by free coordinates. In these
coordinates the feasible germ is exactly an orthant times a Euclidean
space; inactive inequalities have a strict margin. The derivative of
the inverse chart identifies this orthant cone with $K_*$. The first
condition in \eqref{eq:algebraic-corner-tests}, compactness of its unit
section and continuity give a positive coercivity constant. Its linear
image is a closed polyhedral cone, so the second condition gives a
strictly positive attained distance from $p'(0)$ to that image. The
existing regular-corner theorems now apply, including nonunique
preimages of the first projected point. Matrix ranks are tested by
minors; both conditions in \eqref{eq:algebraic-corner-tests} are
first-order statements in linear equations and inequalities. Nash
branch specifications and their derivatives at algebraic points are
encoded by their polynomial equations and isolating conditions.
\end{proof}

Metric projection, constrained value functions and second-order
sensitivity have a substantial existing theory
\cite{RW98,BS98,BS00}. The regular-corner expansion is used here in
that role, with its proof and normalization retained. The additional
statements are the fixed-format classification and exact membership
rule without its regularity hypotheses, the realization theorem, and
the algebraic verification of the regular regime. The two cubic walls
of Sections~\ref{sec:wall} and \ref{sec:geometric-wall} are explicit
realizations of the linear regime; they are not offered as an exhaustive
list of singular entrance geometries.
